\chapter*{Introduction Générale}
\addcontentsline{toc}{chapter}{Introduction Générale}
\markboth{Introduction Générale}{Introduction Générale}

À l'ère de la transformation numérique, le marché du travail a connu une restructuration majeure de ses canaux de recrutement. Les processus de mise en relation traditionnels basés sur les dépôts physiques de dossiers de candidature ont laissé place à des plateformes digitales omniprésentes. Cependant, cette dématérialisation s'accompagne de nouveaux enjeux critiques, notamment en matière de confiance, de sécurité des données et de vérification d'identité.

Au Maroc, l'écosystème du recrutement en ligne fait face à des dérives croissantes. La prolifération de faux recruteurs et d'offres d'emploi frauduleuses sur des canaux informels (réseaux sociaux généralistes, groupes de messagerie) engendre une méfiance partagée : les candidats craignent le vol d'informations personnelles ou les arnaques financières, tandis que les entreprises légitimes souffrent d'une perte d'attractivité et d'usurpations de marque. Par ailleurs, la rigidité des portails d'emploi classiques, qui se limitent souvent à un simple archivage unidirectionnel de CV, ne répond plus aux attentes d'interactivité et de valorisation en temps réel des profils professionnels.

C'est dans ce contexte que s'inscrit le développement de \textbf{ManageraHub}, une plateforme web intégrée combinant recrutement professionnel classique, outils de suivi avancés (Applicant Tracking System - ATS) et réseau social dynamique. La particularité fondamentale de ManageraHub réside dans l'introduction d'une vérification administrative des recruteurs basée sur des identifiants fiscaux et commerciaux officiels (ICE, RC), établissant ainsi un périmètre de confiance garanti.

Ce rapport détaille la démarche méthodologique adoptée pour la concrétisation de ce projet. Il est articulé autour de trois chapitres principaux :
\begin{itemize}
    \item \textbf{Chapitre 1 : Contexte général et présentation du projet}, traitant de l'étude du marché de l'emploi en ligne au Maroc, de l'analyse comparative des solutions existantes, de la critique de l'existant, de la définition de la solution ainsi que de la planification opérationnelle sous le processus unifié \textbf{2TUP}.
    \item \textbf{Chapitre 2 : Analyse et Conception}, détaillant le recueil exhaustif des besoins fonctionnels et non-fonctionnels, les spécifications ergonomiques et graphiques, la modélisation UML complète ainsi que la conception relationnelle de la base de données.
    \item \textbf{Chapitre 3 : Interfaces et Réalisation}, exposant l'environnement de développement, l'implémentation de la logique métier sous Django, la mise en œuvre de la charte graphique et la présentation des écrans finaux via des conteneurs de captures d'écran, ainsi que la validation de la sécurité et des tests.
\end{itemize}
\newpage
